\documentclass[11pt]{amsart}
\usepackage[utf8]{inputenc}
\usepackage[T1]{fontenc}
\usepackage{lmodern}
\usepackage{amsmath,amssymb,amsthm}
\usepackage{mathtools}
\usepackage{booktabs}
\usepackage{array}
\usepackage{enumitem}
\usepackage{xcolor}
\usepackage[margin=1.05in]{geometry}
\usepackage[hidelinks]{hyperref}
\usepackage{url}
\usepackage{listings}
\lstset{basicstyle=\ttfamily\small,breaklines=true,columns=fullflexible,keepspaces=true}

\setlist{leftmargin=*,itemsep=2pt,topsep=3pt}
\allowdisplaybreaks

\newtheorem{theorem}{Theorem}[section]
\newtheorem{lemma}[theorem]{Lemma}
\newtheorem{conjecture}[theorem]{Conjecture}
\theoremstyle{definition}
\newtheorem{remark}[theorem]{Remark}

\newcommand{\col}{\mathrm{col}}
\renewcommand{\UrlFont}{\ttfamily\small}
\makeatletter\g@addto@macro{\UrlBreaks}{\do\-\do\_\do\,}\makeatother

\title[Ahmed--Schaal families]{Four infinite families of the
Ahmed--Schaal conjecture on generalized Schur numbers}
\author{Jerry Yong (with computational assistance)}
\date{October 2026}
\thanks{Formal proof in Lean~4 with symbolic SAT certificates.}

\begin{document}

\begin{abstract}
For $k\ge 3$ let $L(k)$ be the equation $x_1+\dots+x_{k-1}=x_k$. The
generalized Schur number $S(3;s,t,u)$ is the least $N$ such that
every $3$-colouring of $[1,N]$ has a solution of $L(s)$ in colour~1,
of $L(t)$ in colour~2, or of $L(u)$ in colour~3, in each case with
every variable in that colour. Ahmed and Schaal (2016) conjectured
that $S(3;s,t,u)=stu-tu-u-1$ whenever $4\le s\le t\le u$. We prove
the conjecture for $(s,t)\in\{(4,4),(4,5),(4,6),(5,5)\}$ and every
$u\ge t$, which gives $S(3;4,4,u)=11u-1$, $S(3;4,5,u)=14u-1$,
$S(3;4,6,u)=17u-1$ and $S(3;5,5,u)=19u-1$. The upper bound for all
large $u$ comes from one propositional formula per pair, in which the
variables are indexed by linear expressions $au+b$ rather than by
integers. This extends the symbolic SAT method of Chang, De~Loera and
Wesley. The finitely many small $u$ are handled directly. The whole
theorem, including a general lower-bound construction, is checked in
Lean~4 against a short statement file written from the definition.
All $98$ SAT certificates are LRAT proofs checked inside Lean and,
separately, by an independent checker.
\end{abstract}

\maketitle

\section{Introduction}

Schur's theorem says that $S(r)$ is finite: every $r$-colouring of
$[1,S(r)]$ has a monochromatic solution of $x+y=z$. Many
generalizations replace $x+y=z$ by the equation $L(k)$ and allow a
different equation in each colour. A solution may repeat values. For
three colours and parameters $3\le s\le t\le u$, Ahmed and Schaal
\cite{AS16} proved the general lower bound $S(3;s,t,u)\ge stu-tu-u-1$
and computed small values. They made two conjectures. Conjecture~1
concerns $s=3$ and was settled in 2026 \cite{CDW26}. Conjecture~2 is
the subject of this paper.

\begin{conjecture}[Ahmed--Schaal \cite{AS16}]
For $4\le s\le t\le u$, $S(3;s,t,u)=stu-tu-u-1$.
\end{conjecture}

Before this work the conjecture had been verified only for finitely
many triples by direct computer search. The tables in
\cite{AS16,CDW26} list $16$ triples with $s\ge4$, all with $u\le10$.
The best general upper bound, from \cite{CDW26}, comes from the
multicolour Ramsey number $R_3(s,t,u)$. It grows like $u^{s+t-3}$ up
to logarithmic factors, so it is far from linear in $u$. In the
diagonal case the analogous identity $S(3;k,k,k)=k^3-k^2-k-1$ was
proved for all $k$ by Boza et al.\ \cite{Boza}, and later by a
symbolic SAT computation \cite{CDLW}.

\begin{theorem}
Let $M=st-t-1$. For $(s,t)\in\{(4,4),(4,5),(4,6),(5,5)\}$ and every
integer $u\ge t$, $S(3;s,t,u)=Mu-1$. Thus Conjecture~2 holds for
these four infinite families.
\end{theorem}

Note that $stu-tu-u-1=(st-t-1)u-1=Mu-1$, and that $M$ is exactly the
two-colour number $S(2;s,t)$. The proof has three parts:
\begin{enumerate}
\item a lower-bound construction for every $u$ (Section~2);
\item one symbolic formula per pair that gives the upper bound for
all $u\ge u_0$ (Section~3);
\item direct certificates for $t\le u<u_0$ (Section~4).
\end{enumerate}
Section~5 describes the Lean formalization and Section~6 the
computations.

\section{The lower bound}

The general lower bound is already in \cite{AS16}. We give the
following explicit form because it is the one formalized.

\begin{lemma}\label{lem:lower}
Let $\chi_2$ be a $2$-colouring of $[1,M-1]$ with no colour-1
solution of $L(s)$ and no colour-2 solution of $L(t)$. For $u\ge t$,
define a colouring $\chi$ of $[1,Mu-2]$ by $\chi(x)=\chi_2(x)$ for
$x\le M-1$, $\chi(x)=3$ for $M\le x\le (u-1)M-1$, and
$\chi(x)=\chi_2(x-(u-1)M+1)$ for $x\ge (u-1)M$. Then $\chi$ has none
of the three forbidden monochromatic solutions. Hence
$S(3;s,t,u)\ge Mu-1$.
\end{lemma}

\begin{proof}
Colour~3 is the middle block. Any $u-1$ summands from it have sum at
least $(u-1)M$, which lies outside the block. Now take colour
$i\in\{1,2\}$, with parameter $k\le u$, and a solution with $k-1$
summands from the bottom block $[1,M-1]$ and the top block. There are
four cases.

If two or more summands are in the top block, the sum is at least
$2(u-1)M>Mu-2$.

If no summand is in the top block, the sum is at most
$(k-1)(M-1)<(u-1)M$, so the target is in the bottom block. Then the
solution is a $\chi_2$-solution, which is impossible.

If exactly one summand is in the top block, write it as
$(u-1)M-1+y$. The sum is at least $(u-1)M$, so the target is in the
top block too, say $(u-1)M-1+z$. Then $z=y+(\text{the other
summands})$, which is a $\chi_2$-solution of $L(k)$ in colour~$i$.
This is impossible.
\end{proof}

For each pair a suitable $\chi_2$ was found by SAT. It is checked by
evaluation in Lean (\texttt{smallColoringValid\_st}). For example,
for $(4,4)$ the colouring of $[1,10]$ is $2\,2\,1\,1\,1\,1\,1\,1\,2\,2$
up to symmetry.

\section{The upper bound for all large \texorpdfstring{$u$}{u}: a symbolic formula}

Fix $(s,t)$ and integers $B,J$. Let $E$ be the set of pairs
$e=(a,b)$ with $0\le a\le M$ and $-B\le b\le B$, excluding $a=0$ with
$b\le0$ and $a=M$ with $b\ge0$. Read $e$ as the linear function
$e(u)=au+b$. Put $u_0=\max(B+1,\,J+1,\,s,\,t,\,4)$. For every $e\in
E$ and $u\ge u_0$, $1\le e(u)\le Mu-1$. The formula $F(s,t,B,J)$ has
the following variables and clauses:

\begin{description}
\item[Colour variables] $\col[e][i]$, for $i\in\{1,2,3\}$, with
exactly one colour per $e$.

\item[Colours 1 and 2] (parameter $k=s$ or $t$). Reachability
variables $R_i[j][e]$, for $1\le j\le k-1$, mean ``$e$ is a sum of
$j$ elements of colour~$i$''. The clauses are
$\col[e][i]\to R_i[1][e]$ and $R_i[j-1][e]\wedge\col[f][i]\to
R_i[j][e+f]$, where addition is coefficient-wise and $e+f\in E$. The
forbidding clause is $\neg R_i[k-1][e]\vee\neg\col[e][i]$.

\item[Colour 3: $c$-padding] The equation $L(u)$ has a growing
number of summands, so no fixed reachability depth suffices. For each
constant $c$ with $1\le c\le\min(B,M)$, the variable $P_c[j][x]$
means ``$x(u)=(u-1-j)c+(\text{a sum of $j$ colour-3 elements})$''.
The clauses are $\col[(0,c)][3]\to P_c[0][(c,-c)]$ and
$P_c[j-1][x]\wedge\col[y][3]\to P_c[j][x-(0,c)+y]$, for $1\le j\le
J$. Each step replaces one of the $u-1$ copies of $c$ by $y$. The
forbidding clause is $\neg P_c[j][x]\vee\neg\col[x][3]$.
\end{description}

\begin{lemma}[Soundness; Lean \texttt{symForces\_of\_check\_forall}]
\label{lem:sound}
If $F(s,t,B,J)$ is unsatisfiable, then $S(3;s,t,U)\le MU-1$ for every
$U\ge u_0$.
\end{lemma}

\begin{proof}
Let $U\ge u_0$, and let $\chi$ be a colouring of $[1,MU-1]$ with no
forbidden solution. Set $\col[e][i]$ true if and only if
$\chi(e(U))=i$. Set $R$ and $P$ to the true representability
predicates at $U$. Every implication clause holds by construction. A
forbidding clause can fail only if it exhibits a monochromatic
solution; for colour~3 that solution has exactly $u-1$ summands,
since $U-1-j\ge0$. Two expressions may take the same value at $U$,
but this is harmless, because the formula only gets weaker when it
treats them as distinct. So $\chi$ induces a satisfying assignment, a
contradiction.
\end{proof}

The formula does not depend on $u$, so a single SAT call settles
infinitely many cases. The idea of indexing variables by polynomials
comes from \cite{CDLW}, who used it for 3-variable Rado equations.
What is new here is the $c$-padding device for equations whose number
of summands grows with the parameter, together with the use of
off-diagonal colour classes.

\begin{table}[h]
\centering
\caption{Certified symbolic formulas (CaDiCaL~2, direct LRAT, then
\texttt{lrattrim}).}
\begin{tabular}{lrrrrrrl}
\toprule
$(s,t)$ & $M$ & $B$ & $J$ & $u_0$ & vars & clauses & solve (s)\ \ trimmed LRAT \\
\midrule
$(4,4)$ & 11 & 20 & 2 & 21 & 18{,}900  &  1{,}996{,}991 & 5.2\ \ 2.6~MB \\
$(4,5)$ & 14 & 30 & 1 & 31 & 32{,}414  &  5{,}218{,}154 & 16.1\ \ 13.3~MB \\
$(4,6)$ & 17 & 30 & 2 & 31 & 64{,}232  & 16{,}004{,}257 & 35.1\ \ 5.5~MB \\
$(5,5)$ & 19 & 30 & 3 & 31 & 100{,}746 & 31{,}176{,}775 & 73.3\ \ 6.0~MB \\
\bottomrule
\end{tabular}
\end{table}

For smaller $B$, and for smaller $J$ at the same $B$, the formula is
satisfiable. We swept $B\in\{6,10,15,20,30\}$ and $J\in\{1,2,3\}$. A
satisfying assignment does not give a real colouring. It only shows
that the relaxation is too weak.

\section{The finitely many small \texorpdfstring{$u$}{u}}

For $t\le u<u_0$ we use the direct encoding on $[1,N]$ with
$N=Mu-1$. It uses the same reachability scheme, with colour~3 at
depth $u-1$. All $94$ instances ($17+26+25+26$) are unsatisfiable.
Their direct-LRAT proofs took $160$~s of solver time in total and
occupy $16.3$~MiB compressed. Each of the $94$ values also has a
satisfying assignment at $N-1$, which an independent checker
validated. That check is redundant given Lemma~\ref{lem:lower}, but
it confirms the encoding.

\section{Lean formalization}

The trusted part is a $33$-line statement file (SHA-256
\texttt{c6ed1e48\dots6f35}). It contains no SAT notions.

\begin{lstlisting}
def MonoSol (N k : Nat) (chi : Nat -> Fin 3) (i : Fin 3) : Prop :=
  exists l : List Nat, l.length = k - 1 /\
    (forall x in l, 1 <= x /\ x <= N /\ chi x = i) /\
    1 <= l.sum /\ l.sum <= N /\ chi l.sum = i
def Forces (N s t u) := forall chi,
    MonoSol N s chi 0 \/ MonoSol N t chi 1 \/ MonoSol N u chi 2
def Avoidable (N s t u) := exists chi,
    not MonoSol N s chi 0 /\ not MonoSol N t chi 1 /\ not MonoSol N u chi 2
def SchurEq (N s t u) := Forces N s t u /\ Avoidable (N - 1) s t u
def ConjFamily (s t) := forall u, t <= u -> SchurEq ((s*t - t - 1)*u - 1) s t u
\end{lstlisting}

The final file proves \texttt{conj\_4\_4 : ConjFamily 4 4}, and
likewise for $(4,5)$, $(4,6)$ and $(5,5)$.

The formulas are generated by Lean functions. Soundness does not rely
on reasoning about the generators. Instead, a pure Boolean validator
decodes every variable and checks that each clause has one of the
shapes in Section~3, including the coefficient identities and
membership in $E$. The proposition ``every clause is valid'' is
established by \texttt{native\_decide}. Soundness is proved once,
generically, for any CNF all of whose clauses are valid. Certificates
are stored gzip-compressed and Base64-encoded and checked by Lean's
built-in LRAT checker (\texttt{Std.Tactic.BVDecide.LRAT.check}). A
faulty decompressor could only make a check fail, never pass. The
finite cases follow the same pattern.

\texttt{\#print axioms} reports only \texttt{propext},
\texttt{Classical.choice} and \texttt{Quot.sound}, plus the auxiliary
axioms that \texttt{native\_decide} generates, one per decided
computation. These cover the validators, the certificate checks and
the $\chi_2$ checks. So the trusted base is Lean's kernel, its
compiler, and the statement file. The full build takes about
$50$~minutes on one core, dominated by the $94$ finite certificate
checks ($23$~min) and the validity check for the $31$M-clause $(5,5)$
formula ($12$~min).

Independently of Lean, every trimmed proof was checked by a
separately written LRAT checker. The $(4,4)$ symbolic formula was
also regenerated from the written design by a second implementation,
in Python, without reading the original generator. It was again
unsatisfiable.

\section{Limitations and outlook}

For $(5,6)$ and $(6,6)$, the $c$-padding relaxation remained
satisfiable at $B=30$ (and for $(5,6)$ also at $B=40$, $J=1$). Larger
runs were limited by memory and disk. We expect richer padding to
extend the method: several padded constants at once, or summands that
are themselves linear in $u$. A two-parameter formula with
expressions polynomial in both $t$ and $u$ could in principle settle
Conjecture~2 for all $t$ at fixed $s$. We did not attempt this.

On novelty, we checked the arXiv literature through October~2026 and
the $55$ works citing \cite{AS16} that Semantic Scholar lists. None
proves Conjecture~2 for an infinite family. We cannot exclude
unpublished work.

\begin{thebibliography}{9}
\bibitem{AS16} T.~Ahmed and D.~J.~Schaal, On generalized Schur
numbers, \emph{Experimental Mathematics} 25 (2016), 213--218.
\bibitem{CDW26} On two conjectures for generalized off-diagonal Schur
numbers, arXiv:2604.11030 (2026).
\bibitem{CDLW} Y.~Chang, J.~A.~De~Loera and W.~J.~Wesley, Rado
numbers and SAT computations, arXiv:2210.03262 (2022).
\bibitem{Boza} L.~Boza, J.~M.~Mar\'in, M.~P.~Revuelta and
M.~I.~Sanz, 3-color Schur numbers, \emph{Discrete Applied
Mathematics} 263 (2019), 59--68.
\bibitem{CFH17} L.~Cruz-Filipe, M.~J.~H.~Heule, W.~A.~Hunt~Jr.,
M.~Kaufmann and P.~Schneider-Kamp, Efficient certified RAT
verification, CADE-26 (2017).
\bibitem{CaDiCaL} A.~Biere et al., CaDiCaL~2.0, CAV (2024).
\end{thebibliography}

\end{document}
